\chapter{Laboratorium 4}
\label{cha:Lab004}
\makeatletter


% *********************************************************************
% w poniższej metryczce należy uzupełnić informacje o:
% - temacie ćwiczenia
% - numer ćwiczenia
% - numer grupy laboratoryjnej
% - numer zespołu
% - datę wykonania ćwiczenia oraz oddania ćwiczenia.
% *********************************************************************

\begin{table}[H]
    \centering
    \renewcommand{\tabularxcolumn}[1]{m{#1}}  % Dostosowanie kolumny do zawartości
    \newcolumntype{C}{>{\centering\arraybackslash}X}
    \begin{tabularx}{\linewidth}{|C|C|}
        \hline
        \multicolumn{2}{|c|}{\makecell{\textbf{\Large{Biometryczne Systemy Zabezpieczeń}} \\ \textbf{Ćwiczenia Laboratoryjne}}} \\ \hline
        \multicolumn{1}{|l|}{Temat}                    &             TRANSFORMATA HOUGHA            \\ \hline
        \multicolumn{1}{|l|}{Ćwiczenie nr:}                    &      04                   \\ \hline
        \multicolumn{1}{|l|}{Autorzy:}                         &   \@author                              \\ \hline
        \multicolumn{1}{|l|}{Grupa laboratoryjna}                       &       02               \\ \hline
        \multicolumn{1}{|l|}{Zespół nr. }                       &    XX                  \\ \hline
        \multicolumn{1}{|l|}{Data wykonana ćwiczenia}         &  27.04.2025     \\ \hline
        \multicolumn{1}{|l|}{Data oddania ćwiczenia  }         &   13.05.2025     \\ \hline
    \end{tabularx}
\end{table}

\section{Cel ćwiczenia}

Celem zajęć jest zapoznanie się z transformacją Hough’a w zadaniu wykrywania linii i okręgów
w testowych obrazach binarnych oraz obrazach tonowanych pozyskiwanych ze skanerów tęczówki
oka.

\section{Wstęp teoretyczny}

\section{Opis wykorzystywanego oprogramowania i narzędzi}

\section{Zadania do wykonania samodzielnego}

\subsection{Zadanie1 }

\subsubsection{Treść}


\subsubsection{Przykładowy program}

\subsubsection{Obrazy}

\section{Wnioski}

\subsection{Zadanie X}

\subsubsection{Treść}

\subsubsection{Przykładowy program}

\subsubsection{Obrazy}

\section{Wnioski}

Na podstawie przeprowadzonych eksperymentów można stwierdzić, że: